\section{Recovery diagnostics}\label{app:recovery-diagnostics}

\subsection{Directional and magnitude assessment}\label{app:directional-diagnostics}

The assessment in Section~\ref{sec:simulation} uses
\ppnum[0]{ir-validation-redraws} article redraws per history, with
stationary-law, reporting-only and bounded shock controls alongside the
common-reporting coverage ladder.
All sixteen mechanism histories satisfy the positivity, baseline-price,
solver and active-domain requirements.
Each history has a common-reporting operating-law target $\beta_h^0$ and a
reported-law target $\beta_h^\lambda$; sampled coefficients can be
compared with either target according to the measurement question.
Sign frequencies and coefficient errors summarize these comparisons.

For an event indicator $S_{hr}$ in history $h$ and article redraw $r$, let
$S_{hr}=0$ when the estimate is undefined.
With $H=16$ histories and $R$ redraws per history, write
$\widehat p_h=R^{-1}\sum_r S_{hr}$.
The equally weighted frequency and stratified Monte Carlo standard error are
\begin{equation}\label{eq:sign-mcse}
  \widehat p=\frac1H\sum_h\widehat p_h,\qquad
  \widehat{\operatorname{se}}(\widehat p)=
  \left[\frac1{H^2(R-1)}\sum_h\widehat p_h(1-\widehat p_h)\right]^{1/2}.
\end{equation}
The histories are enumerated economic paths, so this standard error measures
article-redraw uncertainty within those paths.

Numerical sign classification uses an arithmetic error allowance fixed before
the control results were inspected.
With $\eta=10^{-12}$, axis-transport allowances propagate through squared
distance, approach and dated financial weights to the centered numerator.
A sign is resolved only when that numerator exceeds its propagated
allowance in absolute value; constructed population zeros retain their
exact-zero status.
This arithmetic convention is separate from the economic solver certificates
and preserves a zero financial denominator as an undefined coefficient.

The identity
\begin{equation}\label{eq:recovery-accounting}
  \widehat\beta^\lambda-\beta^0
  =(\widehat\beta^\lambda-\beta^\lambda)+(\beta^\lambda-\beta^0)
\end{equation}
separates article measurement error from observation distortion.
Only the first term vanishes with infinitely many articles under the
maintained sampling law.
For recovery of the operating-law target, normalized RMSE is
\begin{equation}\label{eq:recovery-rmse}
  \mathcal R_s^\lambda=
  \frac{\sqrt{16^{-1}\sum_h E_r[(\widehat\beta_{h,r,s}^\lambda-\beta_h^0)^2]}}
  {\sqrt{16^{-1}\sum_h(\beta_h^0)^2}}.
\end{equation}
Operational recovery requires this normalized RMSE plus twice its stratified
Monte Carlo standard error to be at most 25\%.
The 25\% threshold is a declared accuracy convention applied only to the
tested budgets.

Signed coefficient bias and its Monte Carlo standard error are in
coefficient units; normalized RMSE uses the oracle RMS as its scale.
For squared errors $L_{hr}$ and oracle RMS $B$, the RMSE standard error uses
\begin{equation}\label{eq:recovery-mcse}
  \widehat{\operatorname{se}}(\widehat{\mathcal R})=
  \frac{\sqrt{H^{-2}R^{-1}\sum_h\widehat{\operatorname{Var}}_r(L_{hr})}}
  {2B\sqrt{H^{-1}R^{-1}\sum_{h,r}L_{hr}}}.
\end{equation}
This estimates article-redraw Monte Carlo uncertainty, conditional on the
enumerated histories, rather than uncertainty about another economic
environment.
If an estimate is undefined, the summary records its count and the affected
recovery criterion cannot pass.
The common-reporting population oracle has zero recovery error by definition;
the changing-reporting population error is calculated exactly.

These summaries are equal-history performance averages; the theoretical expected
moment does not require each finite-history ratio to be positive.
The ratio of expected moments differs from the average of history-specific ratios.
Budget comparisons pair each history's nested redraws rather than treating
budgets or dyads as independent replications.
All attempted estimates remain in their eligible rate denominators;
execution failures, failed economic premises and undefined statistics are
counted separately.
For zero population coefficients, absolute bias, MAE and RMSE remain
available, while sign recovery and normalized RMSE are unavailable.

% AUTO-GENERATED from saved Paper 6 results; do not edit.
\begin{table}[htbp]
\centering
\caption{Directional measurement under fixed economic histories. Targets distinguish common-reporting operating coefficients from reported-law coefficients, including the zero operating target in reporting-only cells. Rates are percentages (stratified Monte Carlo SE in percentage points), not test power. RMSE is in coefficient units against each history's matching target. Appendix~\ref{app:directional-diagnostics} defines resolved nonzero oracles, eligible rate denominators and undefined-attempt accounting. The main table's mean target is an arithmetic mean of history coefficients, not a ratio of expected moments. Each row retains \ppnum[0]{ir-validation-histories} histories and \ppnum[0]{ir-validation-redraws} redraws per history. All attempts enter the rate denominators. No undefined estimates, failed premises or domain departures. Correct sign is unavailable for zero or unresolved targets.}
\label{tab:directional-validation-details}
\begin{tabular}{@{}l@{\enspace}l@{\enspace}r@{\enspace}r@{\enspace}r@{}}
\toprule
Cell & Target & Correct & Positive & RMSE \\
\midrule
\ppvalue{ir-validation-0-label} & \ppvalue{ir-validation-0-target} & \ppnum[1]{ir-validation-0-correct-pct}(\ppnum[2]{ir-validation-0-correct-se-pct}) & \ppnum[1]{ir-validation-0-positive-pct}(\ppnum[2]{ir-validation-0-positive-se-pct}) & \ppnum[3]{ir-validation-0-rmse} \\
\ppvalue{ir-validation-1-label} & \ppvalue{ir-validation-1-target} & \ppnum[1]{ir-validation-1-correct-pct}(\ppnum[2]{ir-validation-1-correct-se-pct}) & \ppnum[1]{ir-validation-1-positive-pct}(\ppnum[2]{ir-validation-1-positive-se-pct}) & \ppnum[3]{ir-validation-1-rmse} \\
\ppvalue{ir-validation-2-label} & \ppvalue{ir-validation-2-target} & \ppnum[1]{ir-validation-2-correct-pct}(\ppnum[2]{ir-validation-2-correct-se-pct}) & \ppnum[1]{ir-validation-2-positive-pct}(\ppnum[2]{ir-validation-2-positive-se-pct}) & \ppnum[3]{ir-validation-2-rmse} \\
\ppvalue{ir-validation-3-label} & \ppvalue{ir-validation-3-target} & \ppnum[1]{ir-validation-3-correct-pct}(\ppnum[2]{ir-validation-3-correct-se-pct}) & \ppnum[1]{ir-validation-3-positive-pct}(\ppnum[2]{ir-validation-3-positive-se-pct}) & \ppnum[3]{ir-validation-3-rmse} \\
\ppvalue{ir-validation-4-label} & \ppvalue{ir-validation-4-target} & \ppnum[1]{ir-validation-4-correct-pct}(\ppnum[2]{ir-validation-4-correct-se-pct}) & \ppnum[1]{ir-validation-4-positive-pct}(\ppnum[2]{ir-validation-4-positive-se-pct}) & \ppnum[3]{ir-validation-4-rmse} \\
\ppvalue{ir-validation-5-label} & \ppvalue{ir-validation-5-target} & \ppnum[1]{ir-validation-5-correct-pct}(\ppnum[2]{ir-validation-5-correct-se-pct}) & \ppnum[1]{ir-validation-5-positive-pct}(\ppnum[2]{ir-validation-5-positive-se-pct}) & \ppnum[3]{ir-validation-5-rmse} \\
\ppvalue{ir-validation-6-label} & \ppvalue{ir-validation-6-target} & \ppmissing & \ppnum[1]{ir-validation-6-positive-pct}(\ppnum[2]{ir-validation-6-positive-se-pct}) & \ppnum[3]{ir-validation-6-rmse} \\
\ppvalue{ir-validation-7-label} & \ppvalue{ir-validation-7-target} & \ppmissing & \ppnum[1]{ir-validation-7-positive-pct}(\ppnum[2]{ir-validation-7-positive-se-pct}) & \ppnum[3]{ir-validation-7-rmse} \\
\ppvalue{ir-validation-8-label} & \ppvalue{ir-validation-8-target} & \ppmissing & \ppnum[1]{ir-validation-8-positive-pct}(\ppnum[2]{ir-validation-8-positive-se-pct}) & \ppnum[3]{ir-validation-8-rmse} \\
\ppvalue{ir-validation-9-label} & \ppvalue{ir-validation-9-target} & \ppnum[1]{ir-validation-9-correct-pct}(\ppnum[2]{ir-validation-9-correct-se-pct}) & \ppnum[1]{ir-validation-9-positive-pct}(\ppnum[2]{ir-validation-9-positive-se-pct}) & \ppnum[3]{ir-validation-9-rmse} \\
\ppvalue{ir-validation-10-label} & \ppvalue{ir-validation-10-target} & \ppmissing & \ppnum[1]{ir-validation-10-positive-pct}(\ppnum[2]{ir-validation-10-positive-se-pct}) & \ppnum[3]{ir-validation-10-rmse} \\
\ppvalue{ir-validation-11-label} & \ppvalue{ir-validation-11-target} & \ppnum[1]{ir-validation-11-correct-pct}(\ppnum[2]{ir-validation-11-correct-se-pct}) & \ppnum[1]{ir-validation-11-positive-pct}(\ppnum[2]{ir-validation-11-positive-se-pct}) & \ppnum[3]{ir-validation-11-rmse} \\
\ppvalue{ir-validation-12-label} & \ppvalue{ir-validation-12-target} & \ppnum[1]{ir-validation-12-correct-pct}(\ppnum[2]{ir-validation-12-correct-se-pct}) & \ppnum[1]{ir-validation-12-positive-pct}(\ppnum[2]{ir-validation-12-positive-se-pct}) & \ppnum[3]{ir-validation-12-rmse} \\
\bottomrule
\end{tabular}
\end{table}


To examine the local approximation, we halve and double the baseline
opportunity innovation at observed article coverage.
All sixteen histories remain in the economic domain and retain positive
tangent and finite-capacity oracles.
Correct-sign recovery is
\ppnum[1]{ir-validation-half-correct-pct}\% and
\ppnum[1]{ir-validation-double-correct-pct}\%, respectively.
The financial denominator scales approximately with the squared shock
amplitude while the population ratio changes little, so this variation
chiefly changes signal relative to article noise.

\subsection{Historical recovery and multi-activity
comparisons}\label{app:historical-recovery}

The 128-redraw study provides the common- and changing-reporting comparison
and the multi-activity dimension checks alongside the current controls.
Its retained magnitude results also document the 25\% recovery benchmark.
Its positive-estimate diagnostic uses
$S_{hr}=\mathbf1\{\widehat\beta_{h,r,s}^\lambda>0\}$ for defined
estimates and zero otherwise, with $R=128$ in \eqref{eq:sign-mcse}.
It precedes the resolved-sign convention used by the current assessment.
Table~\ref{tab:directional-recovery} reports these historical frequencies.

The historical 128-redraw study checks positive opportunities, baseline prices
and active capacities, the solver-distance bound, and agreement of the
operating and affine active domains for every canonical history.
The largest relative solver bound is
\ppnume[2]{ir-maximum-solver-bound}; the minimum canonical capacity and
baseline price are \ppnum[6]{ir-minimum-capacity} and
\ppnum[6]{ir-minimum-baseline-price}.
These checks yield \ppnum[0]{ir-failures} failed histories and
\ppnum[0]{ir-active-domain-violations} active-domain violations.

% AUTO-GENERATED from saved Paper 6 results; do not edit.
\begin{table}[htbp]
\centering
\caption{Historical positive-estimate frequencies from finite article panels, using the original literal $\widehat\beta>0$ convention. All sixteen common-reporting population coefficients are positive. Frequencies average 128 redraws within each history and weight histories equally. SE is the stratified Monte Carlo standard error in percentage points. Zero and undefined estimates do not count as positive; undefined draws remain in the denominator.}
\label{tab:directional-recovery}
\begin{tabular}{llrrr}
\toprule
Reporting & \shortstack{Article-count\\multiplier} & Positive (\%) & SE (pp) & Undefined \\
\midrule
Common & 1 & \ppnum[1]{ir-common-1-positive-pct} & \ppnum[2]{ir-common-1-positive-se-pct} & \ppnum[0]{ir-common-1-sign-undefined} \\
Common & 4 & \ppnum[1]{ir-common-4-positive-pct} & \ppnum[2]{ir-common-4-positive-se-pct} & \ppnum[0]{ir-common-4-sign-undefined} \\
Common & 16 & \ppnum[1]{ir-common-16-positive-pct} & \ppnum[2]{ir-common-16-positive-se-pct} & \ppnum[0]{ir-common-16-sign-undefined} \\
Changing & 1 & \ppnum[1]{ir-changing-1-positive-pct} & \ppnum[2]{ir-changing-1-positive-se-pct} & \ppnum[0]{ir-changing-1-sign-undefined} \\
Changing & 4 & \ppnum[1]{ir-changing-4-positive-pct} & \ppnum[2]{ir-changing-4-positive-se-pct} & \ppnum[0]{ir-changing-4-sign-undefined} \\
Changing & 16 & \ppnum[1]{ir-changing-16-positive-pct} & \ppnum[2]{ir-changing-16-positive-se-pct} & \ppnum[0]{ir-changing-16-sign-undefined} \\
\bottomrule
\end{tabular}
\end{table}


\begin{figure}[tb]
  \centering
  \ppfigure{0.95}{information_recovery}
  \caption{Article information and reporting distortion in the fixed recovery
    experiment.
    Left: normalized recovery RMSE at observed article coverage and count
    multipliers four and sixteen, with bars of two stratified Monte Carlo
    standard errors; population-law points show the exact
    infinite-article error.
    The horizontal line is the 25\% operational convention; the left ordinate
    is logarithmic above that convention and linear below it.
    Right: exact population distortion across the specified log-intensity
    amplitudes, including a fixed-capacity reporting-only observation
    counterfactual.
    Every calculation retains all sixteen canonical economic histories.
    Financial news is calculated from the local valuation response to permanent
    opportunity innovations along the same operating path.
  Reporting amplitudes are illustrative benchmarks rather than empirical calibrations.}
  \label{fig:information-recovery}
\end{figure}

Table~\ref{tab:recovery-diagnostics} retains the historical signed bias,
normalized RMSE and their Monte Carlo standard errors.

% AUTO-GENERATED from saved Paper 6 results; do not edit.
\begin{table}[htbp]
\centering
\caption{Recovery diagnostics against the finite-capacity common-reporting oracle. Bias and its Monte Carlo standard error (SE) are in coefficient units; RMSE and its SE are normalized by oracle RMS. Undefined estimates remain counted. Histories receive equal weight.}
\label{tab:recovery-diagnostics}
\begin{tabular}{llrrrrr}
\toprule
Reporting & \shortstack{Article-count\\multiplier} & Bias & SE & RMSE & SE & Undefined \\
\midrule
Common & 1 & \ppnum[5]{ir-common-1-bias} & \ppnum[5]{ir-common-1-bias-se} & \ppnum[3]{ir-common-1-rmse} & \ppnum[3]{ir-common-1-rmse-se} & \ppnum[0]{ir-common-1-undefined} \\
Common & 4 & \ppnum[5]{ir-common-4-bias} & \ppnum[5]{ir-common-4-bias-se} & \ppnum[3]{ir-common-4-rmse} & \ppnum[3]{ir-common-4-rmse-se} & \ppnum[0]{ir-common-4-undefined} \\
Common & 16 & \ppnum[5]{ir-common-16-bias} & \ppnum[5]{ir-common-16-bias-se} & \ppnum[3]{ir-common-16-rmse} & \ppnum[3]{ir-common-16-rmse-se} & \ppnum[0]{ir-common-16-undefined} \\
Changing & 1 & \ppnum[5]{ir-changing-1-bias} & \ppnum[5]{ir-changing-1-bias-se} & \ppnum[3]{ir-changing-1-rmse} & \ppnum[3]{ir-changing-1-rmse-se} & \ppnum[0]{ir-changing-1-undefined} \\
Changing & 4 & \ppnum[5]{ir-changing-4-bias} & \ppnum[5]{ir-changing-4-bias-se} & \ppnum[3]{ir-changing-4-rmse} & \ppnum[3]{ir-changing-4-rmse-se} & \ppnum[0]{ir-changing-4-undefined} \\
Changing & 16 & \ppnum[5]{ir-changing-16-bias} & \ppnum[5]{ir-changing-16-bias-se} & \ppnum[3]{ir-changing-16-rmse} & \ppnum[3]{ir-changing-16-rmse-se} & \ppnum[0]{ir-changing-16-undefined} \\
\bottomrule
\end{tabular}
\end{table}


The appendix check varies geometric dimension over 2, 4, 8 and 16 while
holding 100 firms, eight activities, 64 demand nodes and $\kappa=1$ fixed.
It uses common reporting and the observed and sixteenfold article budgets,
with the same dated news histories and monthly discount convention.
These are generic multi-activity economies; the two-opportunity sign
mechanism is not imposed.
Table~\ref{tab:recovery-dimensions} reports their recovery and economic-domain checks.

% AUTO-GENERATED from saved Paper 6 results; do not edit.
\begin{table}[htbp]
\centering
\caption{Bounded dimension check under common reporting. All cells hold 100 firms, eight activities, 64 demand nodes and adjustment weight one fixed. RMSE is normalized by oracle RMS, with stratified Monte Carlo SE. All 8 configurations completed with no undefined estimates.}
\label{tab:recovery-dimensions}
\begin{tabular}{llrr}
\toprule
Dimension & \shortstack{Article-count\\multiplier} & RMSE & SE \\
\midrule
2 & 1 & \ppnum[3]{ir-dimension-2-1-rmse} & \ppnum[3]{ir-dimension-2-1-se} \\
2 & 16 & \ppnum[3]{ir-dimension-2-16-rmse} & \ppnum[3]{ir-dimension-2-16-se} \\
4 & 1 & \ppnum[3]{ir-dimension-4-1-rmse} & \ppnum[3]{ir-dimension-4-1-se} \\
4 & 16 & \ppnum[3]{ir-dimension-4-16-rmse} & \ppnum[3]{ir-dimension-4-16-se} \\
8 & 1 & \ppnum[3]{ir-dimension-8-1-rmse} & \ppnum[3]{ir-dimension-8-1-se} \\
8 & 16 & \ppnum[3]{ir-dimension-8-16-rmse} & \ppnum[3]{ir-dimension-8-16-se} \\
16 & 1 & \ppnum[3]{ir-dimension-16-1-rmse} & \ppnum[3]{ir-dimension-16-1-se} \\
16 & 16 & \ppnum[3]{ir-dimension-16-16-rmse} & \ppnum[3]{ir-dimension-16-16-se} \\
\bottomrule
\end{tabular}
\end{table}

